\section*{Native three-node validation of throughput, placement and resource allocation}
\label{sec:native-threefig-validation}
The reported historical results are restored to main-text Figs.~6--8; their unchanged copies remain in Figs.~S2--S3 and S9. We additionally measure all three questions together in a new native Tide experiment. This is a changed-condition validation: it uses three KVM hosts, one ingress, ten-inference requests, fixed execution estimates and an unshaped LAN. The thesis sources establish three-client aggregation and describe YOLOX in the distributed study, with sliding-window estimation in the later goodput source; they do not establish a mixed YOLOX/FaceNet request graph. The historical distributed function's exact workload implementation remains unrecovered. These Tide-only measurements remain valid for their stated conditions, but are not a same-protocol reproduction and do not independently establish historical cross-platform ratios.

\noindent\textbf{Setup and protocol.}
NJ-k1/NJ-k2/NJ-k3 each expose 16 guest vCPUs. CPUs 0--11 form each application's budget; CPUs 12--15 run the control and measurement processes, with no claim of exclusive physical cores. NJ-k2 hosts a persistent two-CPU Entry. NJ-k3 starts with no function Stub; NJ-k1 has one native cloud Stub. Normal scheduling creates edge Stubs and expires idle edge reservations after 2000~ms; cloud and Entry reservations persist. Each request performs ten actual YOLOX-Nano forward passes on the same JPEG through containerd/PAIX, using ONNX Runtime 1.18.1 with two intra-operation threads. Output shape, finite values, all ten output hashes, request identity, execution location and two-CPU affinity are checked against the frozen inputs and calibration. Network links are unshaped.

The 350-ms response deadline is separate from the 180-s client observation timeout. Independently calibrated node execution estimates are 280/357/178~ms for NJ-k1/NJ-k2/NJ-k3. These fixed light-load estimates are not retuned during evaluation; NJ-k2's 357-ms estimate excludes its function execution under this deadline. Three randomized blocks (seed 2026102001) each include 5--50 requests/s in steps of 5, with 30-s offered windows. Each block also includes a separate 60-s 25-requests/s timeline. Every run starts a fresh owned runtime with cached images and an empty edge function inventory. Sweep baseline/recovery periods are 5/10~s; timeline periods are 10/30~s. Observation additionally includes all client terminal records and at least 3~s after the final one. All 33 runs and 29,250 offers are retained.

\noindent\textbf{Definitions and results.}
Effective completion throughput divides verified completions by the longer of the offered window and the interval from its start to the final terminal client return. Goodput divides verified cohort returns within 350~ms of actual client dispatch by the offered duration. Placement fractions use all offers as their denominator, retaining failures and unknown locations. The upper part of Table~\ref{tab:native-threefig} and Fig.~\ref{fig:native-goodput} use the 30 rate-sweep runs. The table's lower rows retain individual outcomes of all three longer 25-requests/s runs, whose resource traces appear in Fig.~\ref{fig:native-processors}. Intervals are pointwise Student-$t$(2) 95\% intervals across three runs.

\begin{figure}[!htbp]
\centering
\includegraphics[width=\linewidth]{revision-figs/evaluation/native_goodput.pdf}
\caption{Native Tide under the stated 12-vCPU application budgets: completion throughput includes response drain; 350-ms goodput counts timely offered-cohort returns. Points are three-run means with pointwise 95\% $t$(2) intervals; dotted lines mark offered rates.}
\label{fig:native-goodput}
\end{figure}

\begin{table}[!htbp]
\centering\small
\caption{Complete new native measurements: upper rows show three-run sweep means $\pm$ pointwise 95\% interval half-widths; lower rows show individual long-run outcomes. Rates are requests/s; location shares are percentages of all offers.}
\label{tab:native-threefig}
\begin{tabular}{rrrrr}
\toprule
Offered / run & Effective throughput &350-ms goodput & Edge (\%) & Cloud (\%)\\
\midrule
5 & 5.00 $\pm$ 0.01 & 4.94 $\pm$ 0.10 & 98.00 $\pm$ 0.00 & 2.00 $\pm$ 0.00 \\
10 & 9.96 $\pm$ 0.02 & 9.62 $\pm$ 0.72 & 98.11 $\pm$ 1.72 & 1.89 $\pm$ 1.72 \\
15 & 14.92 $\pm$ 0.03 & 14.47 $\pm$ 0.74 & 98.37 $\pm$ 0.32 & 1.63 $\pm$ 0.32 \\
20 & 19.88 $\pm$ 0.04 & 19.49 $\pm$ 0.13 & 98.33 $\pm$ 0.00 & 1.67 $\pm$ 0.00 \\
25 & 24.84 $\pm$ 0.08 & 23.84 $\pm$ 1.54 & 98.18 $\pm$ 0.69 & 1.82 $\pm$ 0.69 \\
30 & 29.63 $\pm$ 0.14 & 7.73 $\pm$ 5.07 & 85.15 $\pm$ 13.47 & 14.85 $\pm$ 13.47 \\
35 & 34.55 $\pm$ 0.11 & 10.28 $\pm$ 3.85 & 76.03 $\pm$ 0.36 & 23.97 $\pm$ 0.36 \\
40 & 39.50 $\pm$ 0.11 & 9.16 $\pm$ 1.58 & 65.11 $\pm$ 6.16 & 34.89 $\pm$ 6.16 \\
45 & 40.97 $\pm$ 3.77 & 2.03 $\pm$ 1.44 & 59.04 $\pm$ 0.49 & 40.96 $\pm$ 0.49 \\
50 & 33.94 $\pm$ 8.32 & 1.04 $\pm$ 1.41 & 51.60 $\pm$ 8.68 & 48.40 $\pm$ 8.68 \\
\midrule
\multicolumn{5}{l}{Separate 60-s timeline runs (individual values)}\\
25 (B1, 60s) & 24.84 & 2.80 & 95.40 & 4.60 \\
25 (B2, 60s) & 24.91 & 24.58 & 99.07 & 0.93 \\
25 (B3, 60s) & 24.90 & 19.52 & 98.87 & 1.13 \\
\bottomrule
\end{tabular}
\end{table}

At 25 requests/s, the sweep means are 24.84 completed requests/s and 23.84 timely requests/s. At 50 requests/s, completion throughput reaches only 33.94 requests/s and goodput falls to 1.04 requests/s. Dynamic allocation therefore does not ensure deadline compliance beyond the configured workload's capacity. The longer 25-requests/s runs also show substantial deadline variation; all three remain in the table and plots.

\noindent\textbf{Resource and destination witnesses.}
A locked read-only processor-manager snapshot is sampled from each node every 100~ms; normal heartbeat reservations are independently sampled every 25~ms. Figure~\ref{fig:native-processors} preserves all three timeline runs, summing the per-node observations in each sampling round. Solid lines count committed logical CPUs, including native borrowing; dashed lines count native reserved holdings. These are allocation records, not CPU utilization. Idle Stub expiry releases reservations but can retain pooled model processes. Entry and cloud reservations each have a persistent two-CPU baseline. Missing or invalid observations remain gaps. Figure~\ref{fig:native-placement} links native receiver identity to validated results and same-host executor creation/dispatch records.

\begin{figure}[!htbp]
\centering
\includegraphics[width=\linewidth]{revision-figs/evaluation/native_processors.pdf}
\caption{All three native 25-requests/s timeline runs. Shading marks the 60-s offered window. Each node has a 12-vCPU application budget. Solid/dashed curves sum the three nodes' committed CPUs (including borrowing) and reserved CPUs in each sampling round. The origin's two-CPU Entry and cloud's two-CPU reservation persist after edge expiry; pooled model memory can also persist.}
\label{fig:native-processors}
\end{figure}

\begin{figure}[!htbp]
\centering
\includegraphics[width=.56\linewidth]{revision-figs/evaluation/native_placement.pdf}
\caption{New native execution destinations in the 30 rate-sweep runs. All offered requests remain in the denominator. Points show three-run means and pointwise 95\% $t$(2) intervals; zero-only unknown/failure categories are reported in the text.}
\label{fig:native-placement}
\end{figure}

Across all 33 runs, 29,250 returns pass inference and execution-witness validation; 0 offers fail validation or time out and 0 successful returns have unknown execution locations. Verified returns account for 292,500 actual forward passes. Successful completion, deadline compliance and dynamic allocation are measured separately. The individual long-run rows retain between-run deadline variation even when allocation traces appear similar. The frozen isolated-service estimates can understate service under concurrent execution; deadline compliance depends on the accuracy of these execution estimates.


\noindent\textbf{Evidence and retained interruptions.}
Independent analysis recomputes each request outcome, denominator and location from raw traces, and checks archive hashes, worker/source identities, CPU accounting, executor dispatch joins and owned cleanup. A zero-offer preflight exposed an array-sizing error when four control CPUs and twelve application CPUs were visible to different parts of initialization. An isolated three-line dimension repair precedes all scored runs and leaves scheduling rules unchanged. One first-block45-requests/s run completed all 1,350 client returns before its final cleanup check encountered a process-exit/read race. The original failure record and guarded recovery are retained; that measured cohort is included unchanged, and only unstarted cases continue under the repaired cleanup controller. Per-run cleanup checks restore settings and preserve shared services/tasks. No completed treatment is replaced, and no old/new measurements are combined to recover the original comparative curves.

The original all-instrumentation audit remains flagged: in the first 60-s timeline run, the final valid heartbeat-to-collection-end interval is 527.91~ms, exceeding the frozen 500-ms criterion. This boundary occurs at 90.02--90.55~s after offered start, after the final client return at 60.39~s. Across the matrix, 934 stale/invalid heartbeat observations are retained; 61 belong to this run (10 baseline, 27 offered-window and 24 recovery observations). No threshold is relaxed. Separate metric-specific verification rederives request outcomes and locations using native creation/dispatch logs without heartbeat lifetime attribution; the independently checked 100-ms local processor snapshots cover all nodes throughout this run (largest gap 114.90~ms). Figures~\ref{fig:native-goodput}--\ref{fig:native-placement} use these specifically qualified request and processor records, not filled heartbeat traces. Heartbeats are scheduler inputs and their staleness is not dismissed as harmless; the low first-run goodput remains part of the reported outcomes. The original audit flag, raw observations and separate qualification are retained with their hashes.
